\documentclass[12pt]{article}
\usepackage[utf8]{inputenc}
\usepackage[margin=1in]{geometry}   % 设置页边距
\usepackage{amsmath, amssymb}       % 数学符号与公式
\usepackage{enumitem}               % 自定义列表
% \usepackage{unicode-math} % removed: requires XeLaTeX/LuaLaTeX; keep file pdflatex-friendly
\usepackage[english]{babel}         % 设置整个文档为英文
\usepackage{graphicx}


% 标题与作者信息
\title{Solution for Chapter 5}
\author{Zimo Guo}
\date{}

\begin{document}

\maketitle

% 使用 \noindent 去掉段落缩进，让题号更清晰
\noindent \textbf{5.1}
\begin{enumerate}[label=(\alph*), left=0pt]
    \item 
    When $p = 2$, 
    \begin{align*}
        \mathcal{L}(\mathbf{W}, b, \alpha, \beta, \xi) &= \frac{1}{2}||\mathbf{w}||^2 + 
            C\sum_{i=1}^{m} \xi_{i}^{2} - \sum_{i=1}^{m} \alpha_{i}[y_{i}(\mathbf{w}\mathbf{x}_{i} + b) - 1 + 
            \xi_{i}] - \sum_{i=1}^{m} \beta_{i}\xi_{i}
    \end{align*}

    Setting the gradient of the Lagrangian with respect to the $\mathbf{w}, b, \xi_{i}$ is to zero, the results we got 
    is same as (5.26), (5.27), (5.29), (5.30). (5.28) is replaced by $2C\xi_{i}=\alpha_{i} + \beta_{i}$.

    \begin{align*}
        \inf \mathcal{L} &= \frac{1}{2}||\sum_{i=1}^{m}\alpha_{i}y_{i}\mathbf{x}_{i}||^{2} + 
            C\sum_{i=1}^{m}(\frac{\alpha_{i}+\beta_{i}}{2C})^{2} - \sum_{i=1}^{m}\alpha_{i}y_{i}b + 
            \sum_{i=1}^{m}\alpha_{i} - \sum_{i=1}^{m}\alpha_{i}\xi_{i} - \sum_{i=1}^{m}\beta_{i}\xi_{i} - 
            \sum_{i=1}^{m}\alpha_{i}y_{i}\mathbf{w}\mathbf{x} \\
        &= \sum_{i=1}^{m}\alpha_{i} - \frac{1}{2}\sum_{i=1}^{m}\alpha_{i}\alpha_{j}y_{i}y_{j}\mathbf{x}_{i}\mathbf{x}_{j}
            - \frac{1}{4}\sum_{i=1}^{m}\frac{(\alpha_{i} + \beta_{i})^{2}}{C}
    \end{align*}

    \item 
    Yes, this problem is still convex because the first two items are convex, and $-\frac{1}{4}\sum_{i=1}^{m}\frac{(\alpha_{i} + \beta_{i})^{2}}{C}$ is 
    concave but we can use convec techniques to solve it.
\end{enumerate}

\vspace{1em}  % 题目之间的垂直间距

\noindent \textbf{5.2}

\vspace{1em}  % 题目之间的垂直间距

\noindent \textbf{5.3}

Add weight to Lagrangian, geo 

\begin{align*}
    \mathcal{L}(\mathbf{W}, b, \alpha, \beta, \xi) &= \frac{1}{2}||\mathbf{w}||^2 + 
        C\sum_{i=1}^{m} \xi_{i}p_{i} - \sum_{i=1}^{m} \alpha_{i}[y_{i}(\mathbf{w}\mathbf{x}_{i} + b) - 1 + 
        \xi_{i}] - \sum_{i=1}^{m} \beta_{i}\xi_{i}
\end{align*}

For gradient, (5.26)(5.27)(5.29)(5.30) is same as regular case. (5.28) is replaced with $Cp_{i} - 
\alpha_{i} - \beta_{i}$. Put them into Lagrangian, get the same result with regular case.

\begin{align*}
    \inf \mathcal{L} &= \sum_{i=1}^{m}\alpha_{i} - \frac{1}{2}\sum_{i=1}^{m}\alpha_{i}\alpha_{j}y_{i}y_{j}\mathbf{x}_{i}\mathbf{x}_{j}
\end{align*}

The constraint condition is $\sum_{i=1}^{m}\alpha_{i}y_{i}=0 \wedge 0 \le \alpha_{i}, \beta_{i} \le Cp_{i}$

\vspace{1em}  % 题目之间的垂直间距

\noindent \textbf{5.4}

\begin{enumerate}[label=(\alph*), left=0pt]
    \item 
    \begin{align*}
        Equation (5.33) &= \alpha_{1}+\alpha_{2}+\sum_{i=3}^{m}\alpha_{i} \\
        &- \frac{1}{2}[\alpha_{1}\alpha_{1}y_{1}y_{1}(\textbf{x}_{1})\textbf{x}_{1} 
        + \alpha_{1}\alpha_{2}y_{1}y_{2}(\textbf{x}_{1}\textbf{x}_{2}) 
        + \sum_{i=3}^{m}\alpha_{1}\alpha_{i}y_{1}y_{i}(\textbf{x}_{1}\textbf{x}_{i}) \\
        &+ \alpha_{2}\alpha_{1}y_{2}y_{1}(\textbf{x}_{2}\textbf{x}_{1}) 
        + \alpha_{2}\alpha_{2}y_{2}y_{2}(\textbf{x}_{2}\textbf{x}_{2}) 
        + \sum_{i=3}^{m}\alpha_{2}\alpha_{i}y_{2}y_{i}(\textbf{x}_{2}\textbf{x}_{i})] \\
        &= \alpha_{1}+\alpha_{2}-\frac{1}{2}K_{11}\alpha_{1}^{2}-\frac{1}{2}K_{22}\alpha_{2}^{2}-sK_{12}\alpha_{1}\alpha_{2}-y_{1}\alpha_{1}v_{1}-y_{2}\alpha_{2}v_{2}
    \end{align*}

    For constraint condition, \\
    \begin{align*}
        0 = \sum_{i=1}^{m}\alpha_{i}y_{i} &= \alpha_{1}y_{1}+\alpha_{2}y_{2}+\sum_{i=3}^{m}\alpha_{i}y_{i} \\
        &= \alpha_{1}y_{1}^{2}+\alpha_{2}y_{2}y_{1}+y_{1}\sum_{i=3}^{m}\alpha_{i}y_{i} \\
        &= \alpha_{1} + s\alpha_{2} + \gamma = 0
    \end{align*}
\end{enumerate}

\vspace{1em}  % 题目之间的垂直间距

\noindent \textbf{5.6}

\begin{enumerate}[label=(\alph*), left=0pt]
    \item First item in original optimization problem is $||\mathbf{w}||^{2}$, 
        \begin{align*}
            ||\mathbf{w}||^{2} &= ||\sum_{i=1}^{m}\alpha_{i}y_{i}\mathbf{x}_{i}||^{2} \\
            &= (\alpha_{1}y_{1}\mathbf{x}_{1}+\alpha_{2}y_{2}\mathbf{x}_{2}+\dots+\alpha_{m}y_{m}\mathbf{x}_{m})(\alpha_{1}y_{1}\mathbf{x}_{1}+\alpha_{2}y_{2}\mathbf{x}_{2}+\dots+\alpha_{m}y_{m}\mathbf{x}_{m})
        \end{align*}
        
        To make above equation equal to $\sum_{i=1}^{m}\alpha_{i}^{2}$, following should holds 
        \begin{align*}
            \mathbf{x}_{i}\mathbf{x}_{j} = 0, i \ne j
        \end{align*}.

    \item 
    \begin{align*}
        &L=\frac{1}{2}\sum_{i=1}^{m}\alpha_{i}^{2}+C\sum_{i=1}^{m}\xi_{i}-\sum_{i=1}^{m}\beta_{i}\cdot[y_{i}(\sum_{j=1}^{m}\alpha_{j}y_{j}\textbf{x}_{i}\cdot \textbf{x}_{j}+b)-1+\xi_{i}]-\sum_{i=1}^{n}\gamma_{i}\xi_{i}-\sum_{i=1}^{m}\varepsilon_{i}\alpha_{i} \\
        &\text{KKT condition:} \\
        & \nabla_{\alpha_{i}}L=\alpha_{i} - \sum_{j=1}^{m}\beta_{j}y_{j}y_{i}\textbf{x}_{i}\textbf{x}_{j}-\varepsilon_{i}=0 \\
        &\nabla_{b}L=\sum_{i=1}^{m}\beta_{i}\cdot y_{i}=0 \\
        &\nabla_{\xi_{i}}L=C-\beta_{i}-\gamma_{i}=0 \\
        &\beta_{i}[y_{i}(\sum_{j=1}^{m}\alpha_{j}y_{j}\textbf{x}_{i}\textbf{x}_{j}+b)-1+\xi_{i}]=0 \\
        &\gamma_{i}\cdot\xi_{i}=0 \\
        &\varepsilon_{i}\cdot\alpha_{i}=0
    \end{align*}

    \item 
    \begin{align*}
        &L=\sum_{i=1}^{m}||\alpha_{i}||+C\sum_{i=1}^{m}\xi_{i}-\sum_{i=1}^{m}\beta_{i}\cdot[y_{i}(\sum_{j=1}^{m}\alpha_{j}y_{j}\textbf{x}_{i}\cdot \textbf{x}_{j}+b)-1+\xi_{i}]-\sum_{i=1}^{n}\gamma_{i}\xi_{i}-\sum_{i=1}^{m}\varepsilon_{i}\alpha_{i} \\
        &\text{KKT condition:} \\
        & \nabla_{\alpha_{i}}L=1 - \sum_{j=1}^{m}\beta_{j}y_{j}y_{i}\textbf{x}_{i}\textbf{x}_{j}-\varepsilon_{i}=0 \\
        &\nabla_{b}L=\sum_{i=1}^{m}\beta_{i}\cdot y_{i}=0 \\
        &\nabla_{\xi_{i}}L=C-\beta_{i}-\gamma_{i}=0 \\
        &\beta_{i}[y_{i}(\sum_{j=1}^{m}\alpha_{j}y_{j}\textbf{x}_{i}\textbf{x}_{j}+b)-1+\xi_{i}]=0 \\
        &\gamma_{i}\cdot\xi_{i}=0 \\
        &\varepsilon_{i}\cdot\alpha_{i}=0
    \end{align*}
\end{enumerate}


\vspace{1em}  % 题目之间的垂直间距

\noindent \textbf{5.7}

\begin{enumerate}[label=(\alph*), left=0pt]
    \item 
    Because ${\textbf{x}_{1}, \dots, \textbf{x}_{d}}$ can be shattered, there always exists
    $\mathbf{w}$ which can holds $\mathbf{w}\mathbf{x}_{x}=y_{i}$. 

    \begin{align*}
        1 &\le |\mathbf{w}\mathbf{x}_{i}| \\
        &= |\mathbf{w}y_{i}\mathbf{x}_{i}| \\
        &= \mathbf{w}y_{i}\mathbf{x}_{i}
    \end{align*}

    Add above for $i = 1, \dots, d$, get

    \begin{align*}
        d &\le \sum_{i=1}^{d}\mathbf{w}y_{i}\mathbf{x}_{i} \\
        &\le \mathbf{w}\sum_{i=1}^{d}y_{i}\mathbf{x}_{i} \\
        d &\le ||\mathbf{w}|| ||\sum_{i=1}^{m}y_{i}\mathbf{x}_{i}|| \\
        &\le \Lambda||\sum_{i=1}^{m}y_{i}\mathbf{x}_{i}|| 
    \end{align*}

    \item 
    According to Jesen's inequality, $E(x)^{2} \le E(x^{2})$, so $E(x) \le \sqrt{E(x^{2})}$.

    \begin{align*}
        d &\le \Lambda E[||\sum_{i=1}^{m}y_{i}\mathbf{x}_{i}||] \\
        &\le \Lambda \sqrt{E[||\sum_{i=1}^{m}y_{i}\mathbf{x}_{i}||^{2}]} \\
        &= \Lambda \sqrt{\sum_{i=1}^{d}||\mathbf{x}_{i}||^{2}}
    \end{align*}

    \item 
    \begin{align*}
        d &\le \Lambda \sqrt{\sum_{i=1}^{d}||\mathbf{x}_{i}||^{2}} \\
        &\le \Lambda \sqrt{dr^{2}} \\
    \end{align*}

    Solve the inequality, get $d \le r^{2}\Lambda^{2}$.
\end{enumerate}


\end{document}